% Generated draft: requires mathematical and source review.
\documentclass[11pt]{article}
\usepackage[margin=25mm]{geometry}
\usepackage{fontspec}
\setmainfont{Latin Modern Roman}
\usepackage{amsmath,amssymb,mathtools,unicode-math}
\setmathfont{Latin Modern Math}
\usepackage{enumitem,needspace,xurl,hyperref}
\setlength{\parindent}{0pt}
\setlength{\parskip}{5pt}
\setlength{\emergencystretch}{3em}
\newcommand{\N}{\mathbb{N}}
\newcommand{\Z}{\mathbb{Z}}
\newcommand{\Q}{\mathbb{Q}}
\newcommand{\R}{\mathbb{R}}
\newcommand{\C}{\mathbb{C}}
\newcommand{\F}{\mathbb{F}}
\newcommand{\A}{\mathbb{A}}
\newcommand{\PP}{\mathbb{P}}
\newcommand{\E}{\mathbb{E}}
\newcommand{\eps}{\varepsilon}
\newcommand{\dd}{\,\mathrm{d}}
\begin{document}
\section*{UnsolvedMath ID: 20001974}
\textbf{Source problem code:} catalogue entry
\subsection{Definitions and mathematical statement}
T. Holm's Question \(9.3\) asks whether the middle-perversity intersection cohomology of a singular K\"ahler quotient has a natural ring structure. ``Natural'' is not specified. A precise candidate requirement is independence of a chosen small resolution: for a proper small resolution \(\pi:Y\to X\) from a smooth variety, use the normalized intersection-cohomology identification \(T_\pi:IH^*(X;\Q)\cong H^*(Y;\Q)\) and transport cup product by
\[
 a\star_\pi b=T_\pi^{-1}(T_\pi(a)\smile T_\pi(b)).
\]
Small means that, for every \(r>0\), the locus whose fiber dimension is at least \(r\) has complex codimension greater than \(2r\). The candidate asks whether these products agree for different small resolutions under their natural additive identifications, rather than merely whether the resulting rings are abstractly isomorphic.

Comment \(9.4\), E. Lerman, attributes small-resolution constructions and induced products for circle reductions to E. Lerman--S. Tolman [LT00], noting their special circle-action setting. Comment \(9.5\), R. Sjamaar, cites Kiem--Woolf's example with two small resolutions and concludes that a natural product appears obstructed. The supplied aid makes the sense of that discrepancy explicit.

For the Kiem--Woolf \(\C^*\)-action on \(\mathbb P^4\) with weights \((0,1,1,-1,-1)\), take the singular quotient and its two small resolutions in the cited construction. In the aid's common rational intersection-cohomology model,
\[
 V=\operatorname{span}_\Q\{1,h,u,h^2,hu,h^3\},\qquad
 |h|=|u|=2,
\]
where the degree-four symbol \(hu\) is a nonzero basis element, the two naturally transported products satisfy
\[
 u\star_+u=-h^2+2hu,\qquad
 u\star_-u=-h^2-2hu,
\]
\[
 (u\star_+u)-(u\star_-u)=4hu\ne0.
\]
Thus cup product cannot be transported from an arbitrary small resolution in a resolution-independent way under these identifications. At the same time, the aid reports that the two abstract graded rings are isomorphic by \(h\mapsto h\), \(u\mapsto-u\). The source phrase ``distinct cohomology rings'' must therefore be understood as distinct naturally transported products here, not as abstract nonisomorphism.

The example rules out that particular resolution-independent construction. It does not prove that no product can be chosen on the vector space or that every conceivable enriched, chamber-dependent or differently natural construction is impossible. The exact quotient linearization and definitions of the common-model generators are not supplied in this record and require the cited example. The appended heading ``\(10\) Computations over \(\Z\)'' belongs to the next section and supplies no further question.
\subsection{Short English statement}
Two small resolutions can induce different products on the same naturally identified intersection-cohomology space. The aid gives an explicit discrepancy \(4hu\) in the Kiem--Woolf example, while the two rings remain abstractly isomorphic. This obstructs resolution-independent transport of cup product, rather than every possible notion of a natural product.
\subsection{Sources}
\begin{itemize}
\item Local UnsolvedMath registry, permanent ID \(20001974\), source code \texttt{catalogue entry}; canonical collection \texttt{aim-geometry-notes.json}, notes[311].
\item AIM workshop ``Moment maps and surjectivity in various geometries,'' Question \(9.3\), T. Holm; Comment \(9.4\), E. Lerman; Comment \(9.5\), R. Sjamaar: \url{https://aimath.org/WWN/momentmaps/momentmaps.pdf}. Page number \(15\) and the Section \(10\) heading are extraction artifacts or spillover. No PDF was opened in this session.
\item E. Lerman and S. Tolman [LT00], small resolutions for \(S^1\)-reductions, without full title or publication data in the record.
\item Kiem and Woolf, ``Intersection cohomology of symplectic quotients by circle actions.'' The original comment says ``to be published in J. London Math. Soc.'' and drops the initial I in extraction; the aid describes the example as published. No publication year or complete citation is supplied.
\item Supplied research attempt \(1\), \texttt{partial\_result}, run status \texttt{partially\_solved}, rigor label ``proved,'' novelty confidence ``low.'' Registry classification \texttt{PARTIAL-PROGRESS}, updated 14 August 2026. Original remarks and literature fields are empty; no proposer-field value or proposed year is supplied.
\end{itemize}
\end{document}
